% mode: ini-etex
% Aimed at §1067: the penalty between lines.  \interlinepenalty, \clubpenalty, \widowpenalty, \displaywidowpenalty and
% \brokenpenalty, and the e-TeX arrays \interlinepenalties, \clubpenalties, \widowpenalties, \displaywidowpenalties that
% replace them, shorter and longer than the paragraph, with the clamp of r at the last element; d=true from a display;
% a sum of zero gives no node; one- and two-line paragraphs, where club and widow meet.
\input prelude
\hyphenchar\tenrm=`\-
\tracingparagraphs=0 \parindent=0pt \pretolerance=-1 \tolerance=10000 \hbadness=10000 \hfuzz=16383pt
\showboxbreadth=1000 \showboxdepth=1 \hsize=100pt \baselineskip=12pt
\abovedisplayskip=0pt \belowdisplayskip=0pt \abovedisplayshortskip=0pt \belowdisplayshortskip=0pt
\def\L#1{\penalty-10000 \hbox to#1pt{}}
\def\Two{\noindent\hbox to 1pt{}\L2}
\def\Three{\Two\L3}
\def\Six{\Three\L4\L5\L6}
\def\T#1{\setbox1\vbox{#1\par}\showbox1 }
\T\Two \T\Three \T\Six \T{\noindent\hbox{}}
\interlinepenalty=1 \T\Two \T\Three \T\Six
\clubpenalty=20 \T\Two \T\Three \T\Six
\widowpenalty=300 \T\Two \T\Three \T\Six
\interlinepenalty=-20 \T\Six
\clubpenalty=0 \widowpenalty=0 \interlinepenalty=0
% a sum that is zero
\interlinepenalty=5 \clubpenalty=-5 \T\Three
\interlinepenalty=5 \clubpenalty=0 \widowpenalty=-5 \T\Three
\interlinepenalty=0 \clubpenalty=0 \widowpenalty=0
% \brokenpenalty only after a discretionary break
\brokenpenalty=4000
\T{\noindent\hbox to 95pt{}\discretionary{}{}{}\hbox to 95pt{}\hskip 5pt\hbox to 95pt{}\discretionary{a}{b}{}\hbox to 90pt{}}
\brokenpenalty=0
% the arrays, short, exact and long; \interlinepenalties counts lines of the paragraph, the others from its ends
\interlinepenalties=1 7 \T\Six
\interlinepenalties=2 7 8 \T\Six
\interlinepenalties=7 1 2 3 4 5 6 7 \T\Six
\interlinepenalties=0 \interlinepenalty=99 \T\Six
\interlinepenalty=0
\clubpenalties=1 7 \T\Six
\clubpenalties=3 7 8 9 \T\Six
\clubpenalties=9 1 2 3 4 5 6 7 8 9 \T\Six
\clubpenalties=0 \clubpenalty=99 \T\Six
\clubpenalty=0
\widowpenalties=1 7 \T\Six
\widowpenalties=3 7 8 9 \T\Six
\widowpenalties=9 1 2 3 4 5 6 7 8 9 \T\Six
\widowpenalties=0 \widowpenalty=99 \T\Six
\widowpenalty=0
% all at once, with values that cancel out in a line
\interlinepenalties=3 1 2 3 \clubpenalties=2 10 20 \widowpenalties=2 100 200 \T\Six
\interlinepenalties=2 -30 0 \clubpenalties=2 30 -2 \T\Six
\interlinepenalties=0 \clubpenalties=0 \widowpenalties=0
% with \prevgraf: \interlinepenalties goes by the line in the page, \clubpenalties by cur_line-prev_graf
\interlinepenalties=7 1 2 3 4 5 6 7 \clubpenalties=4 10 20 30 40
\T{\noindent\prevgraf=2 \Three\L4}
\interlinepenalties=0 \clubpenalties=0
% from a display, line_break(true)
\displaywidowpenalty=5000 \widowpenalty=70
\def\D{$$\vrule width 10pt height 5pt$$}
\T{\Three\D\Three}
\displaywidowpenalties=3 1000 2000 3000 \widowpenalties=2 11 22
\T{\Six\D\Six}
\displaywidowpenalties=0 \T{\Six\D\Six}
\widowpenalties=0 \T{\Six\D\Six}
\displaywidowpenalty=0 \widowpenalty=0 \clubpenalty=33 \T{\Three\D\Three}
\clubpenalty=0
% very large penalties add up past 2^31 on the oracle: only the sign can be compared, so no box is shown here
\interlinepenalty=2000000000 \clubpenalty=2000000000 \widowpenalty=2000000000
\def\Sign{\setbox2\lastbox \unskip \count3=\ifnum\lastpenalty>0 1\else\ifnum\lastpenalty<0 -1\else0\fi\fi\relax
  \immediate\write-1{sign of the penalty \the\count3}\unpenalty}
\setbox1\vbox{\Three\par \Sign \Sign}
\interlinepenalty=-2000000000 \clubpenalty=-2000000000 \widowpenalty=-2000000000
\setbox1\vbox{\Three\par \Sign \Sign}
\end
